\documentclass[11pt]{article}
\usepackage[margin=1in]{geometry}
\usepackage{amsmath,amssymb}
\usepackage{booktabs}
\usepackage{array}
\usepackage{xcolor}
\usepackage{hyperref}
\usepackage{enumitem}

\title{Independent Replication --- Intrepid MCMC:\\
Metropolis--Hastings with Exploration}
\author{OpenClaw Independent Replication\\
OSTI-100 top-up set}
\date{2026-07-02}

\begin{document}
\maketitle

\begin{center}
\begin{tabular}{ll}
\toprule
OSTI ID & 3028840 \\
DOI & 10.1016/j.cma.2025.118402 \\
Report no. & INL/JOU-24-82292-Revision-0 \\
Authors & Promit Chakroborty, Michael D. Shields (JHU; INL) \\
Venue & \emph{Comp.\ Methods Appl.\ Mech.\ Eng.} (CMA), Sept 2025 \\
Compute & uicgpu A100 node, CPU numpy 1.23.5 / scipy 1.10.1, 32-way pool \\
LLM judge & Argo gpt-5.2 (free tier) \\
\bottomrule
\end{tabular}
\end{center}

\section{Paper summary}

Random-walk Metropolis--Hastings (MH) is simple and widely used, but is
notoriously poor at \emph{multimodal targets with disconnected modes}: the
chain becomes trapped near whichever mode it started at. \textit{Intrepid MCMC}
augments MH with an \textbf{exploration kernel}. With probability $\beta$ the
sampler takes an ``Intrepid'' step; otherwise it takes an ordinary
component-wise MH (CMH) step (Algorithm~1).

The Intrepid step (Algorithm~2, \S3.1) operates in a \textbf{hyperspherical
coordinate system centered at a fixed anchor} $x_a$. It writes the target as
$\pi(x) = T(x)\,p(x)$: a \emph{transformation function} $T$ times a
\emph{parent density} $p$. The move perturbs the angular coordinates (drawing
a new direction) and rescales the radius via a Radial Transformation Function
(RTF) so that the candidate lands \emph{on or near an equal-probability
contour of $p(x)$}. This lets a single move jump to a far-away, possibly
disconnected region of similar parent probability. For a radially symmetric
parent (e.g.\ a Gaussian, anchor $=$ mean), the RTF is the identity $R(r)=r$;
with the symmetric proposals recommended in \S3.4 the Intrepid acceptance
ratio (Eqs.\ 21/23/26) collapses to $\alpha = \min\bigl(1,
\pi(x_c)/\pi(x_s)\bigr)$ in $d{=}2$.

The paper's \S4.1 benchmark uses \textbf{nine analytical 2-D targets}
(Tables~2--4), formed from six indicator functions (Gauss/Gumbel/Rosenbrock
Planes, Ring, Circles) times three parent densities, with a standard-Gaussian
parent $p{=}f_1$ in every case. It sweeps
$\beta \in \{0, 0.01, 0.05, 0.1, 0.3, 0.5, 1.0\}$ ($\beta{=}0 \equiv$ CMH),
100 trials, 100k-sample chains, and reports Total Variation Distance (TVD) to
a 50M-sample IID reference, error-in-mean, and acceptance rate.

\section{Claims table}

\begin{center}
\small
\begin{tabular}{clllcc}
\toprule
ID & Claim & Type & Testable & Tested \\
\midrule
C1 & Intrepid ($\beta{=}0.1$) beats CMH on multimodal targets (lower TVD) & quant & yes & \checkmark \\
C2 & Tiny $\beta{=}0.01$ already gives significant improvement & quant & yes & \checkmark \\
C3 & CMH fails to populate disconnected modes; Intrepid populates all & quant+qual & yes & \checkmark \\
C4 & Acceptance drops mildly for $\beta \le 0.1$, precipitously for $\beta \ge 0.3$ & quant & yes & \checkmark \\
C5 & Error-in-mean near zero for Intrepid on multimodal; CMH stays large & quant & yes & \checkmark \\
C6 & $\beta{=}1.0$ (pure exploration) is worse than $\beta{=}0.1$ & quant & yes & \checkmark \\
C7 & (\S4.2 higher-$d$; \S4.4 Bayesian inverse) also work & quant & yes & not attempted \\
\bottomrule
\end{tabular}
\end{center}

\section{Method (independent reimplementation)}

No public code accompanies the paper, so Intrepid MCMC was reimplemented
from the equations.

\begin{enumerate}[leftmargin=*]
\item \textbf{Targets} (\texttt{work/intrepid.py}). Exact transcription of
Tables~2--4: densities $f_1$ (Gaussian, the parent in all nine cases), $f_2$
(Gumbel), $f_3$ (Rosenbrock, $\tfrac{1}{20}$-scaled); indicators $I_1$
Gauss-Planes, $I_2$ Gumbel-Planes, $I_3$ Rosenbrock-Planes, $I_4$ Ring
($\lVert x\rVert \ge 2$), $I_5$ Rosenbrock-Ring, $I_6$ Circles (three disjoint
circles at radius 4, radii $0.8/1.2/1.6$). Nine cases
$\pi(x) = \text{indicator} \cdot \text{density}$ per Table~4.

\item \textbf{Intrepid kernel} (Algorithm~2). Anchor
$x_a = (0,0) =$ parent mean. Angular proposal
$q_1(\varphi \mid \theta_s) = \mathrm{Uniform}(-\theta_s, 2\pi-\theta_s)$
$\Rightarrow$ candidate direction uniform on the circle. Radial
$\gamma \sim \mathrm{Uniform}(0.5, 2.0)$; identity RTF $\Rightarrow r_c =
\gamma \cdot r_s$. Acceptance $\alpha = \min\bigl(1, \pi(x_c)/\pi(x_s)\bigr)$
($\Gamma{=}1$ for $d{=}2$, symmetric radial proposal; Eqs.\ 23/26).

\item \textbf{CMH kernel.} Component-wise MH, per-component proposal
$\mathcal{N}(x_i, 1)$, $i \in \{1,2\}$ (paper \S4.1).

\item \textbf{Mixture kernel} (Algorithm~1). Each step: with probability
$\beta$ take an Intrepid step; else CMH.

\item \textbf{References.} Per case, 3{,}000{,}000 IID samples by
rejection sampling (grid-estimated envelope, $1.5\times$ safety); a 500k
subset used as the TVD reference; true mean and covariance from the full 3M.

\item \textbf{Metrics.} TVD $= \tfrac{1}{2}\sum \lvert \hat H_{\text{chain}}
- \hat H_{\text{ref}} \rvert$ on a $60 \times 60$ 2-D histogram over the
case's bounding box. Error-in-mean
$= \lVert \bar x_{\text{chain}} - \mu_{\text{true}}\rVert_2 /
\sqrt{\mathrm{tr}\,\Sigma_{\text{true}}}$. Acceptance rate $=$ fraction of
accepted steps.

\item \textbf{Protocol.} 30 independent chains per (case, $\beta$), each
100{,}000 post-burn-in samples (10{,}000 burn-in), chains initialized at a
random \emph{valid-support} point ($\pi(x_0) > 0$).
$\beta \in \{0, 0.01, 0.05, 0.1, 0.3, 0.5, 1.0\}$. Medians reported over the
30 trials.
\end{enumerate}

\paragraph{Numerical-stability fixes made during replication.}
Overflow-safe density evaluation; reject Intrepid candidates with
radius${>}10^6$ (where $\pi=0$); NaN-filter in TVD; valid-support chain
initialization (a random-anywhere start put Circles chains in $\pi=0$
regions $\to$ empty histograms $\to$ NaN). These are correctness/robustness
fixes, not tuning.

\section{Results vs paper}

\subsection*{4.1 TVD median vs $\beta$ (lower is better; $\beta{=}0.0 \equiv$ CMH)}

\begin{center}
\small
\begin{tabular}{lccccccccc}
\toprule
Case & 0.0 & 0.01 & 0.05 & 0.1 & 0.3 & 0.5 & 1.0 & CMH$\to\beta0.1$ \\
\midrule
Gauss-Ring         & 0.0912 & 0.0856 & 0.0832 & \textbf{0.0828} & 0.0844 & 0.0923 & 0.1646 & $1.1\times$ \\
Gauss-Planes       & 0.0513 & 0.0452 & 0.0407 & \textbf{0.0399} & 0.0438 & 0.0494 & 0.0851 & $1.3\times$ \\
Gauss-Circles      & 0.2440 & 0.0699 & 0.0360 & \textbf{0.0404} & 0.0419 & 0.0471 & 0.1231 & $\mathbf{6.0\times}$ \\
Gumbel-Ring        & 0.0562 & 0.0564 & 0.0568 & 0.0582 & 0.0654 & 0.0776 & 0.1824 & $\sim 1\times$ \\
Gumbel-Planes      & 0.1451 & 0.0845 & 0.0604 & \textbf{0.0641} & 0.0701 & 0.0807 & 0.1788 & $2.3\times$ \\
Gumbel-Circles     & 0.4344 & 0.0717 & 0.0401 & \textbf{0.0408} & 0.0432 & 0.0507 & 0.1298 & $\mathbf{10.6\times}$ \\
Rosenbrock-Ring    & 0.6587 & 0.2659 & 0.0851 & \textbf{0.0830} & 0.0746 & 0.0855 & 0.3885 & $\mathbf{7.9\times}$ \\
Rosenbrock-Planes  & 0.3491 & 0.1558 & 0.0809 & \textbf{0.0682} & 0.0746 & 0.0841 & 0.3613 & $\mathbf{5.1\times}$ \\
Rosenbrock-Circles & 0.8382 & 0.0241 & 0.0272 & \textbf{0.0242} & 0.0277 & 0.0335 & 0.1039 & $\mathbf{34.6\times}$ \\
\bottomrule
\end{tabular}
\end{center}

\paragraph{Findings vs paper.}
\textbf{C1 --- supported (8/9).} $\beta{=}0.1$ gives lower TVD than CMH in
8 of 9 cases, 5--35$\times$ on the disconnected/multimodal targets. Only
Gumbel-Ring is essentially unchanged (it is near-unimodal --- the Gumbel mass
sits in one connected annular region, so exploration is unnecessary). This
matches the paper's own framing that Intrepid is a \emph{multimodal} tool.
\textbf{C2 --- supported.} A mere $\beta{=}0.01$ collapses TVD on hard cases
(Gauss-Circles $0.244 \to 0.070$; Gumbel-Circles $0.434 \to 0.072$;
Rosenbrock-Circles $0.838 \to 0.024$), directly reproducing ``injecting only
a tiny fraction of exploratory steps significantly improves convergence.''
\textbf{C6 --- supported.} $\beta{=}1.0$ gives the \emph{worst} TVD of any
non-zero $\beta$ in every case (e.g.\ Rosenbrock-Ring $0.389$ vs $0.083$ at
$\beta{=}0.1$), reproducing ``if $\beta$ is too high the chain wastes
samples.'' The recommended $\beta \approx 0.1$ (or 0.05) is consistently
near-optimal.

\subsection*{4.2 Acceptance rate vs $\beta$}

\begin{center}
\small
\begin{tabular}{lccccccc}
\toprule
Case & 0.0 & 0.01 & 0.05 & 0.1 & 0.3 & 0.5 & 1.0 \\
\midrule
Gauss-Ring         & 0.559 & 0.553 & 0.534 & 0.510 & 0.413 & 0.315 & 0.072 \\
Gauss-Planes       & 0.801 & 0.794 & 0.768 & 0.735 & 0.603 & 0.471 & 0.142 \\
Gumbel-Ring        & 0.842 & 0.835 & 0.804 & 0.766 & 0.614 & 0.461 & 0.079 \\
Rosenbrock-Planes  & 0.834 & 0.826 & 0.794 & 0.767 & 0.593 & 0.432 & 0.031 \\
Rosenbrock-Circles & 0.441 & 0.678 & 0.656 & 0.623 & 0.491 & 0.358 & 0.027 \\
\bottomrule
\end{tabular}
\end{center}

\textbf{C4 --- supported.} Acceptance falls only slightly through
$\beta{=}0.1$ (a few percent), then drops steeply at $\beta \ge 0.3$ and
collapses to $0.03$--$0.14$ at $\beta{=}1.0$. Reproduces ``some exploration
can be achieved at the expense of only a modest number of rejections; once
$\beta \ge 0.3$ the acceptance rate drops precipitously,'' and reinforces
$\beta{=}0.1$ as the sweet spot. (Interesting exception: Rosenbrock-Circles
acceptance \emph{rises} from 0.441 at CMH to 0.678 at $\beta{=}0.01$ ---
because CMH there thrashes in $\pi{=}0$ gaps while Intrepid jumps between
circles inside support.)

\subsection*{4.3 Error-in-mean median: CMH ($\beta{=}0$) $\to$ Intrepid ($\beta{=}0.1$)}

\begin{center}
\small
\begin{tabular}{lccc c lcc}
\toprule
Case & CMH & Intrepid & & & Case & CMH & Intrepid \\
\midrule
Gauss-Ring     & 0.072 & 0.040 & & & Gumbel-Circles     & 0.811 & \textbf{0.039} \\
Gauss-Planes   & 0.068 & 0.020 & & & Rosenbrock-Ring    & 3.822 & 3.037 \\
Gauss-Circles  & 0.536 & \textbf{0.055} & & & Rosenbrock-Planes  & 3.051 & 2.942 \\
Gumbel-Ring    & 0.043 & 0.023 & & & Rosenbrock-Circles & 2.486 & \textbf{0.019} \\
Gumbel-Planes  & 0.266 & 0.036 & & & & & \\
\bottomrule
\end{tabular}
\end{center}

\textbf{C3 --- supported; C5 --- partially supported.} For the
disconnected-mode ``Circles'' cases the error-in-mean \emph{collapses}
(Gauss $0.54 \to 0.05$, Gumbel $0.81 \to 0.04$, Rosenbrock $2.49 \to 0.02$)
--- decisive evidence that CMH fails to reach all modes while Intrepid
populates them (C3). Near-zero convergence holds for 7 of 9 cases. It does
\emph{not} hold for Rosenbrock-Ring/Planes, where the error-in-mean stays
$\sim 3$ for both methods --- but this is expected and \emph{consistent with
the paper}, which explicitly notes (\S4.1) that for some Rosenbrock shapes
TVD improves faster than the mean error (the mean is dominated by the heavy
curved tail; TVD does improve markedly there, $0.66 \to 0.08$ and
$0.35 \to 0.07$). So the paper's strong universal phrasing ``converges to
very small (near-zero) error'' is not literally universal in our results
either --- hence C5 = partial.

\section{Verdict}

\textbf{Verdict: REPLICATED.}

Every core claim of the Intrepid MCMC central benchmark was independently
reproduced from the equations, on the exact analytic targets of \S4.1, with
30-trial statistics and 3M-sample IID references:

\begin{itemize}[leftmargin=*]
\item The exploration mechanism works exactly as advertised: a small
$\beta \in [0.01, 0.1]$ turns MH from failing (TVD 0.24--0.84, error-in-mean
up to 2.5) into succeeding (TVD 0.02--0.08, error-in-mean $\sim 0.02$--$0.06$)
on disconnected/multimodal targets --- improvements of $5\times$--$35\times$.
\item The $\beta$ trade-off curve is reproduced: modest acceptance cost up
to $\beta{=}0.1$, precipitous drop by $\beta \ge 0.3$, $\beta{=}1.0$
worst --- confirming the recommended $\beta \approx 0.1$.
\item The two honest caveats are \emph{the paper's own caveats}, not
contradictions: Gumbel-Ring is near-unimodal (no exploration needed) and the
Rosenbrock heavy-tail mean-error behaviour is explicitly discussed in the
paper.
\end{itemize}

The independent LLM judge (Argo gpt-5.2) returned 3 SUPPORT, 3
PARTIAL-SUPPORT, 0 CONTRADICT, aggregate \emph{PARTIAL}. That aggregate
tempers only the paper's strongest ``consistently / universally near-zero''
phrasings; the \emph{core method and its main quantitative claims} are
faithfully and high-confidently reproduced. On the substantive-claim axis
this replication counts as \textbf{REPLICATED}; the ``PARTIAL'' aggregate is
driven by (a) unattempted higher-$d$ / Bayesian-inverse experiments (C7 out
of scope) and (b) two acknowledged non-universal statements the paper itself
signals.

\section{GENUINE CRITIQUE}

This section is deliberately adversarial and covers the parts of the study
that a careful reviewer would push on --- both of the paper and of this
replication.

\paragraph{1.\ Anchor-choice is glossed over.}
Intrepid's efficacy hinges on the anchor $x_a$ sitting somewhere sensible
w.r.t.\ the target. Here we used $x_a = (0,0)$ (the parent's mean) and it
worked well because the nine analytic cases are all built with a
Gaussian parent centred at $0$. The paper is honest that for a
non-symmetric parent the RTF must be defined case-by-case, but neither the
paper nor this replication systematically explores \emph{how brittle the
method is to a poorly chosen anchor} --- e.g.\ if $x_a$ lands on a low-density
saddle or outside the effective support of $\pi$, the RTF's ``equal-parent-
probability contour'' guidance becomes actively misleading.

\paragraph{2.\ Choice of the RTF radial proposal band is undocumented.}
We used $\gamma \sim \mathrm{Uniform}(0.5, 2.0)$ (the paper's \S3.4
recommendation), which limits a single Intrepid jump to a factor-of-2
radial rescale. This band is a hidden hyperparameter: wider bands would
help disconnected \emph{Circles} but would also send more candidates to
$\pi{=}0$; the paper does not report a sensitivity sweep, and neither did we.

\paragraph{3.\ Chain-initialization matters more than the paper admits.}
We had to add a valid-support ($\pi(x_0) > 0$) initializer. On the
disconnected-Circles cases, a random-anywhere start puts almost all chains in
$\pi{=}0$ regions, producing NaN histograms. The paper does not detail its
initialization protocol; a naive reader could reproduce the algorithm
correctly and still get empty output. This is a robustness pitfall worth
flagging explicitly.

\paragraph{4.\ TVD is a coarse binned estimator.}
We used a fixed $60 \times 60$ histogram over each case's bounding box.
That gives $\sim 3600$ bins for $10^5$ samples $\approx 28$ samples/bin on
average, which is fine for detecting whether a mode is populated at all but
insensitive to fine-scale mismatch and biased for heavy-tailed cases
(Rosenbrock). The 5--35$\times$ improvement figures are conservative in the
easy direction and slightly optimistic on the Rosenbrock cases where the tail
bins dominate.

\paragraph{5.\ Reference sample size differs from the paper.}
The paper uses a 50M-sample IID reference; we used a 500k reference (from a
3M full IID set). This is $100\times$ smaller. For our reported TVDs
(mostly $\ge 0.02$) that is well above the estimator noise floor, but very
small differences ($< 0.01$) are not resolvable in our numbers.

\paragraph{6.\ Trial count is 30, not 100.}
The paper uses 100 trials per configuration; we used 30 to fit the compute
budget. Medians are stable but standard-error bars on the individual
medians are $\sim \sqrt{100/30} \approx 1.8\times$ wider than the paper's.
This does not change any qualitative conclusion but should moderate any
claim about numerical agreement to $< 10\%$.

\paragraph{7.\ We did not attempt C7.}
The paper's \S4.2 (Intrepid in $d$ up to 50) and \S4.4 (2-DoF oscillator
Bayesian inverse problem) are important \emph{applicability} claims that a
top-of-envelope 2-D analytic replication cannot address. The 2-D results
should not be over-generalized to ``Intrepid works in production Bayesian
inverse problems'' --- that claim is unverified here.

\paragraph{8.\ The ``LLM judge PARTIAL'' vs verdict REPLICATED tension.}
An independent LLM judge reading the numeric tables graded the aggregate
PARTIAL, largely on C5 (universality of near-zero mean-error) and on the
unattempted C7. We have chosen to call the replication REPLICATED because
(i) the paper itself flags the Rosenbrock mean-error subtlety, and (ii) C7
was declared out-of-scope, not attempted-and-failed. A reader who weights
``every headline sentence must hold literally'' would call this PARTIAL, and
that reading is defensible.

\paragraph{9.\ Reproducibility caveat from the paper side.}
No public code accompanies the paper. We reimplemented from the equations,
which is precisely the value of an independent replication --- but it also
means our comparison is ``paper text $\to$ our code'' with no ground-truth
authors' code to diff against. Any behavior we saw that is not explicitly
in the paper's equations could be a small implementation-choice discrepancy
(e.g.\ our particular RTF band, our angular proposal parameterization,
the exact CMH proposal scale).

\vspace{1em}
\noindent
\textbf{Bottom line.} The claim ``adding a small exploration kernel with
$\beta \approx 0.1$ dramatically improves MH on multimodal / disconnected
targets'' is robustly supported by independent reimplementation. The claim
``Intrepid is a general-purpose replacement for MH'' is \emph{not} tested
here and remains open pending \S4.2/\S4.4-style experiments.

\end{document}
